\section{前期研究：多算为何收效甚微}

本文的研究问题始于作者在东海赤潮巡测场景下对机载算力分配的初步探索：当自主水下航行器仅携带一张3天前的陈旧卫星叶绿素浓度图下水，需要在“继续机载推演计算”“边走边测”与“停下复测”之间分配有限作业时间时，多算究竟能否带来更好的观测收益\cite{lou2014diurnal,mnr2024bulletin}？在该前期研究中，我们用多年卫星叶绿素数据标定旧图误差，构造48个想象海洋（32个用于搜索航线、16个留出不参与搜索）\cite{beckers2003dineof,gaspari1999correlation}，采用交叉熵方法在机载端搜索巡测航线\cite{deboer2005cem}，并以航行器随时报告的最浓位置与真实最浓处之间的对数叶绿素浓度差的时间平均（遗憾）作为评价指标，考察了推演次数与时间分配策略的作用。

前期研究的第一个核心发现是，在陈旧背景图上大幅增加机载推演计算，只会使规划器在自身搜索过的想象海洋上给出的自评分数持续虚高，而航线在真实海洋中的价值早已停止增长\cite{smith2006curse,vanhasselt2010double}。如图\ref{fig:curse}所示，当每次决策的想象推演次数从64次增加到2000次乃至20万次时，规划器搜索时的自评价值从0.762一路攀升至0.833和0.917；然而，所选航线在真实海洋中的价值仅从0.749升至0.793和0.803，在约2000次推演后便彻底饱和，且在不参与搜索的留出海洋上的独立评判（0.745、0.794、0.802）始终紧贴真实价值。这表明过多的算力并未帮航行器看清真实赤潮的位置，而是被用来过度迎合有限想象样本中的偶然高值噪声。

\begin{figure}[htbp]
\centering
\includegraphics[width=\linewidth]{fig_curse}
\caption{前期东海赤潮巡测研究揭示的机载过度计算现象。(a) 所选航线价值（对数叶绿素浓度）随每次决策想象推演次数（64至20万次）的变化：规划器搜索时的自评从0.762持续升至0.917，而真实价值从0.749升至0.793后在约2000次推演处饱和（20万次时为0.803），留出海洋上的独立评判紧贴真实价值；(b) 估计价值减真实价值的偏差：搜索自评的高估偏差随计算量单调膨胀，留出评判保持无偏，不确定性幅度设错仅引起整体平移。多算只让规划器的答案在旧信息上越来越“精致”，却不能提供区分真实环境的新信息。}
\label{fig:curse}
\end{figure}

前期研究的第二个发现是，在连续平滑的标量场巡测中，无论如何精巧地设计计算与观测的时间分配策略，其所能改变的结局空间都极其有限。实验表明，采用出发前预演选出的用时策略其平均遗憾为0.265，而即使在每个任务结束后事后挑出最好的用时策略，遗憾也只能降到0.243，意味着用时分配本身的最大可改进空间仅有约0.022（经典单步元推理为0.287）。对遗憾来源的拆解进一步显示，当航行器从距离赤潮最浓处中位数达48千米的窗口边缘下水时，知道真相直奔最浓处的遗憾为0.206，换成当天卫星图为0.235；若改在旧图最浓处下水，五种不同用时策略的遗憾均集中在0.126至0.145之间，换成当天卫星图则降至0.098（尝试用前一天光流外推3天反而比直接用旧图误差增大0.018，95\%置信区间为0.012至0.025，说明逐日卫星图里没有可外推的整体移动）。换言之，下水位置决定了约0.13的差距，信息新鲜度决定了约0.03，而计算与观测的用时分配仅占约0.02。

这一前期探索促使作者把问题进一步收紧：既然当多个可能的世界都与现有旧图相容、却要求采取截然不同的行动时，继续计算无法将它们区分开来，那么在强海流绕障这种一旦选错一侧就会付出数小时代价的离散拓扑决策中，我们能否直接检验“是否存在另一个同样合理却让当前选择吃大亏的流场”，并以此决定何时该算、何时该测、何时该保守绕行？这正是本文接下来要解决的核心问题。
